\section*{अध्याय \ref{ch:b3:complex-kinetics} --- जटिल बलगतिकी: शृंखलाएँ, एंज़ाइम और दोलन}

\begin{solution}{exo:b3:complex-kinetics:1}
$\ce{Cl2 -> 2 Cl}$: आरंभन; $\ce{Cl + H2 -> HCl + H}$ और $\ce{H + Cl2 -> HCl + Cl}$:
संचरण; $\ce{2 Cl -> Cl2}$: समापन। वाहक: Cl और H। समग्र (दोनों
संचरण चरणों का योग): $\ce{H2 + Cl2 -> 2 HCl}$।
\end{solution}

\begin{solution}{exo:b3:complex-kinetics:2}
$[\mathrm M]_{1/2} = k_2/k_{-1} = \qty{1.0e-3}{mol/L} = \qty{1.0}{mol/m^3}$; $p = cRT = 1.0 \times 8.314 \times 750 =
\qty{6.2}{kPa}$।
\end{solution}

\begin{solution}{exo:b3:complex-kinetics:3}
$v/V_{\max} = [\mathrm S]/(K_M + [\mathrm S])$: $\frac12$, $\frac34$, $\frac{9}{10}$। $V_{\max}$ के 90\,\% के लिए $[\mathrm S] = 9K_M$ चाहिए।
\end{solution}

\begin{solution}{exo:b3:complex-kinetics:4}
$k_{\mathrm{cat}}/K_M = 100/\num{0.80e-3} = \qty{1.25e5}{L.mol^{-1}.s^{-1}}$: विसरण सीमा से कोई 60\,000 गुना नीचे,
और विशिष्ट मान के निकट।
\end{solution}

\begin{solution}{exo:b3:complex-kinetics:5}
\ce{CH3CO} के लिए स्थायी अवस्था: $k_2[\ce{CH3}][\ce{CH3CHO}] = k_3[\ce{CH3CO}]$। \ce{CH3} के लिए स्थायी अवस्था:
$k_1[\ce{CH3CHO}] - k_2[\ce{CH3}][\ce{CH3CHO}] + k_3[\ce{CH3CO}] - 2k_4[\ce{CH3}]^2 = 0$। जोड़ने पर: $k_1[\ce{CH3CHO}] =
2k_4[\ce{CH3}]^2$, अतः $[\ce{CH3}] = (k_1/2k_4)^{1/2}[\ce{CH3CHO}]^{1/2}$, और $\dd[\ce{CH4}]/\dd t = k_2[\ce{CH3}][\ce{CH3CHO}] =
k_2(k_1/2k_4)^{1/2}[\ce{CH3CHO}]^{3/2}$।
\end{solution}

\begin{solution}{exo:b3:complex-kinetics:6}
आधे रूपांतरण पर $[\ce{H2}] = [\ce{Br2}] = \frac12$, $[\ce{HBr}] = 1$ (आपेक्षिक इकाइयाँ):
$v/v_0 = \frac12 \times (\frac12)^{1/2}/(1 + 0.10 \times 1/0.5) = 0.354/1.2 = 0.29$। निरोधन के बिना यह
0.354 होता: निरोधन दर को और 1.2 से भाग देता है।
\end{solution}

\begin{solution}{exo:b3:complex-kinetics:7}
विस्फोट तब जब $1000p > 6000/p + 10p^2$, अर्थात् $10p^3 - 1000p^2 + 6000 < 0$, जिसके धनात्मक
मूल 2.48 और \qty{99.9}{kPa} हैं: मिश्रण लगभग 2.5 और
\qty{100}{kPa} के बीच विस्फोटित होता है (मॉडल की पहली और दूसरी सीमाएँ)।
\end{solution}

\begin{solution}{exo:b3:complex-kinetics:8}
$1/v = 1/V_{\max} + (K_M/V_{\max})/[\mathrm S]$: $5.0 = 1/V_{\max} + 2.0K_M/V_{\max}$ और $2.5 = 1/V_{\max} + 0.5K_M/V_{\max}$।
घटाने पर: $K_M/V_{\max} = \qty{1.667}{s.mM.\micro M^{-1}}$, फिर $1/V_{\max} = 1.667$: $V_{\max} = \qty{0.60}{\micro M.s^{-1}}$,
$K_M = \qty{1.0}{mM}$।
\end{solution}

\begin{solution}{exo:b3:complex-kinetics:9}
$1/v$ अक्ष पर वही अंतःखंड: स्पर्धी। $K_M^{\mathrm{app}} = 1/0.417 = \qty{2.40}{mM} = 3K_M$, अतः
$1 + 50/K_i = 3$ और $K_i = \qty{25}{\micro M}$।
\end{solution}

\begin{solution}{exo:b3:complex-kinetics:10}
$N = \qty{1.01}{mol/L}$: $t_{1/2} = \ln(1.01/0.010 - 1)/(0.50 \times 1.01) = \ln100/0.505 = \qty{9.1}{s}$।
दर $kx(N - x)$, $x = N/2$ पर अधिकतम है: वही क्षण, S-वक्र का
नति-परिवर्तन बिंदु।
\end{solution}

\begin{solution}{exo:b3:complex-kinetics:11}
अनुरेख $B - 2$, सारणिक 1। $B = 1.5$: $\lambda = -0.25 \pm 0.968\iu$, एक स्थिर केंद्र-सर्पिल (स्थायी
अवस्था की ओर अवमंदित दोलन)। $B = 2.5$: $\lambda = 0.25 \pm 0.968\iu$, एक अस्थिर केंद्र-सर्पिल:
दोलन तब तक बढ़ते हैं जब तक वे \omterm{def:b3:complex-kinetics:oscillating}{सीमा चक्र} तक न पहुँच जाएँ।
\end{solution}

\begin{solution}{exo:b3:complex-kinetics:12}
$1/\tau = \num{1.0e8} \times \num{5.40e-5} + \num{1.0e3} = \qty{6.4e3}{s^{-1}}$, $\tau = \qty{156}{\micro s}$। अनेक
कुल सांद्रताओं पर $\tau$ मापिए: $1/\tau$, $[\mathrm A]_e + [\mathrm B]_e$ के विरुद्ध, ढलान $k_1$ और
अंतःखंड $k_{-1}$ वाली एक सीधी रेखा है।
\end{solution}

\begin{solution}{pb:b3:complex-kinetics:1}
\textbf{1.} आरंभ में $[\mathrm S]$ ज्ञात होता है और उत्पाद, जो संदमन कर सकता है या उलटी अभिक्रिया
कर सकता है, अनुपस्थित होता है।
\textbf{2.} निम्न-$[\mathrm S]$ बिंदु, बड़े $1/[\mathrm S]$ और $1/v$ पर: वे ढलान तय करते हैं, और
व्युत्क्रम उनकी त्रुटियों को बढ़ा देता है।
\textbf{3.} यह हर मापी गई दर को उसी रूप में भार देता है जिसमें वह मापी गई; लाइनवीवर--बर्क रेखा
अभिनत प्राचल देती है।
\textbf{4.} $K_M$: $0.773$, 0.801 से 1.2 मानक त्रुटि नीचे है; $V_{\max}$: $0.491$, 0.502 से 2.2 मानक
त्रुटि नीचे। दोनों कम हैं, जैसा अभिनति से अपेक्षित है।
\textbf{5.} $k_{\mathrm{cat}} = 0.502/0.0050 = \qty{100}{s^{-1}}$।
\textbf{6.} $100/\num{0.801e-3} = \qty{1.25e5}{L.mol^{-1}.s^{-1}}$।
\textbf{7.} विसरण सीमा से कोई 60\,000 गुना नीचे; एक विशिष्ट \omterm{def:b3:complex-kinetics:enzyme}{एंज़ाइम}।
\textbf{8.} स्पर्धी: $V_{\max}$ अपरिवर्तित, आभासी $K_M$, $[\mathrm I]$ के साथ बढ़ता हुआ।
\textbf{9.} $K_M^{\mathrm{app}} = K_M(1 + [\mathrm I]/K_i) = K_M + (K_M/K_i)[\mathrm I]$।
\textbf{10.} $K_i = 0.799/0.0321 = \qty{24.9}{\micro M}$।
\textbf{11.} $24.9 \pm 4.30 \times 0.9 = 24.9 \pm \qty{3.9}{\micro M}$।
\textbf{12.} एंज़ाइम--संदमक संकुल का वियोजन स्थिरांक: वह
संदमक सांद्रता जो मुक्त \omterm{def:b3:complex-kinetics:enzyme}{एंज़ाइम} के आधे भाग को घेर लेती है।
\textbf{13.} $v_i/v_0 = (K_M + [\mathrm S])/(K_M(1 + [\mathrm I]/K_i) + [\mathrm S])$।
\textbf{14.} $2/3 = 0.67$।
\textbf{15.} $2/12 = 0.17$।
\textbf{16.} $11/21 = 0.52$: उच्च क्रियाधार सांद्रता पर क्रियाधार संदमक को
स्पर्धा में पछाड़ देता है।
\textbf{17.} $[\mathrm S] = K_M$ पर $v_i/v_0 = 2/(2 + [\mathrm I]/K_i) = 0.10$ से $[\mathrm I]/K_i = 18$ मिलता है।
\textbf{18.} $18 \times 24.9 = \qty{448}{\micro M}$।
\textbf{19.} पहला चरण तीव्र है: हर \omterm{def:b3:complex-kinetics:enzyme}{एंज़ाइम} अणु शीघ्र एक $\mathrm P_1$ मुक्त करता है
और ऐसिल-एंज़ाइम के रूप में फँस जाता है; उसके बाद $\mathrm P_1$ केवल उतनी तेज़ी से प्रकट होता है जितनी तेज़ी से
धीमा दूसरा चरण \omterm{def:b3:complex-kinetics:enzyme}{एंज़ाइम} को मुक्त करता है।
\textbf{20.} $1.28/2.0 = 0.64 = (k_2/(k_2 + k_3))^2$, अतः $k_2/(k_2 + k_3) = 0.80$।
\textbf{21.} $200/2.0 = \qty{100}{s^{-1}}$, भाग I जितना ही।
\textbf{22.} $k_2 = 0.80 \times 625 = \qty{500}{s^{-1}}$, $k_3 = \qty{125}{s^{-1}}$।
\textbf{23.} $500 \times 125/625 = \qty{100}{s^{-1}}$।
\textbf{24.} \textbf{$[\mathrm I]_{90\,\%} = 18K_i \approx \qty{0.45}{mM}$, $[\mathrm S] = K_M$ पर।}
\end{solution}
